\begin{remark*}
The distance from a point to a set is not a new kind of object. It is a number
defined inside an already chosen ambient metric space.

Mathematically, let $(X,d)$ be a metric space, let $x \in X$, and let
$A \subseteq X$. The distance from $x$ to $A$ is
\[
  \operatorname{dist}(x,A)
  =
  \inf \{ d(x,a) : a \in A \}.
\]
This is an infimum, not necessarily a minimum. If the closest point belongs to
$A$, then the infimum is achieved. If the closest point lies only on the
boundary of $A$, then the infimum may be approached without ever being attained.

In Lean, the ambient space is explicit. The generic declaration is
\[
\begin{gathered}
  \texttt{def distanceSet}\\
  \texttt{\ \ \ \ \{Point : Type u\}}\\
  \texttt{\ \ \ \ (M : MetricSpace Point)}\\
  \texttt{\ \ \ \ (x : Point)}\\
  \texttt{\ \ \ \ (S : Set Point) : Set Real :=}\\
  \texttt{\ \ (fun y : Point => M.distance x y) '' S.}
\end{gathered}
\]
This helper definition names the set
\[
  \{ d(x,y) : y \in S \}.
\]
It is not mathematically necessary: we could write this set inline everywhere.
However, naming it makes the formal statements match the usual hand proof,
where one first writes something like
\[
  D = \{ d(x,y) : y \in S \}
\]
and then proves that the desired distance is the infimum of $D$.

The point-to-set distance is then defined by taking the infimum of that named
set:
\[
\begin{gathered}
  \texttt{noncomputable def distanceToSet}\\
  \texttt{\ \ \ \ \{Point : Type u\}}\\
  \texttt{\ \ \ \ (M : MetricSpace Point)}\\
  \texttt{\ \ \ \ (x : Point)}\\
  \texttt{\ \ \ \ (S : Set Point) : Real :=}\\
  \texttt{\ \ sInf (distanceSet M x S).}
\end{gathered}
\]
Here \texttt{Point} is the carrier type, \texttt{M} is the chosen metric space
structure on that carrier, \texttt{x} is the point, and \texttt{S} is the set.
\texttt{distanceSet M x S} is the set of all real numbers that occur as
\texttt{M.distance x y} for some \texttt{y} in \texttt{S}.

We do not introduce a structure or a named type for this. The point-to-set
distance is a function of the ambient metric, the point, and the set:
\[
  \texttt{distanceToSet M x S}.
\]
Different metrics on the same carrier may give different values, so the metric
space argument must stay visible.
\end{remark*}

\begin{remark*}
The Lean primitive is called \texttt{sInf}, the infimum of a set. In ordinary
mathematical writing we call this simply ``the infimum'' or write
\[
  \inf \{ d(x,a) : a \in A \}.
\]
The formal bridge is the theorem
\[
\begin{gathered}
  \texttt{distanceToSet\_isGLB :}\\
  \texttt{\ \ IsGLB (distanceSet M x S) (distanceToSet M x S).}
\end{gathered}
\]
The hypothesis \texttt{S.Nonempty} is required because the usual real-valued
point-to-set distance is normally discussed for nonempty sets. In this Lean
development, \texttt{distanceToSet} is real-valued. That is convenient for the
current metric-space layer, but it means the empty-set case should not be read
as the usual extended-real convention $\operatorname{dist}(x,\emptyset)=+\infty$.
\end{remark*}

\begin{remark*}
The proof of \texttt{distanceToSet\_isGLB} is the formal version of a standard
hand proof.

Mathematically, write
\[
  D = \{ d(x,y) : y \in S \}.
\]
Since $S$ is nonempty, choose $y_0 \in S$. Then $d(x,y_0) \in D$, so $D$ is
nonempty. Also, every metric distance is nonnegative, so
\[
  0 \le d(x,y)
\]
for every $y \in S$. Therefore $0$ is a lower bound for $D$, and $D$ is bounded
below. Since $\mathbb{R}$ has greatest lower bounds for nonempty bounded-below
sets, $\inf D$ is the greatest lower bound of $D$.

In Lean, this breaks into the same pieces. First, prove that the distance set
is nonempty:
\[
\begin{gathered}
  \texttt{theorem distanceSet\_nonempty}\\
  \texttt{\ \ \ \ \{Point : Type u\}}\\
  \texttt{\ \ \ \ (M : MetricSpace Point)}\\
  \texttt{\ \ \ \ (x : Point)}\\
  \texttt{\ \ \ \ \{S : Set Point\}}\\
  \texttt{\ \ \ \ (set\_nonempty : S.Nonempty) :}\\
  \texttt{\ \ \ \ (distanceSet M x S).Nonempty := by}\\
  \texttt{\ \ sorry.}
\end{gathered}
\]
The proof begins by unpacking \texttt{set\_nonempty} into a point
\texttt{y : Point} and a proof \texttt{point\_in\_set : y in S}. Then
\texttt{M.distance x y} is an explicit element of \texttt{distanceSet M x S}.

Second, prove that the distance set is bounded below:
\[
\begin{gathered}
  \texttt{theorem distanceSet\_bddBelow}\\
  \texttt{\ \ \ \ \{Point : Type u\}}\\
  \texttt{\ \ \ \ (M : MetricSpace Point)}\\
  \texttt{\ \ \ \ (x : Point)}\\
  \texttt{\ \ \ \ (S : Set Point) :}\\
  \texttt{\ \ \ \ BddBelow (distanceSet M x S) := by}\\
  \texttt{\ \ exact M.distance\_image\_bddBelow x S.}
\end{gathered}
\]
The mathematical witness for \texttt{BddBelow} is \texttt{0}. After choosing
that witness, the remaining obligation is to show that every member of
\texttt{distanceSet M x S} is at least \texttt{0}. A member of
\texttt{distanceSet M x S} has the form \texttt{M.distance x y} for some
\texttt{y in S}, and this is exactly the nonnegativity axiom of the metric.
Because \texttt{distanceSet M x S} is defined as the image
\texttt{(fun y : Point => M.distance x y) '' S}, this bounded-below proof can
reuse the general metric-space theorem
\[
  \texttt{M.distance\_image\_bddBelow x S}.
\]

With those two helper facts, the final theorem has the shape
\[
\begin{gathered}
  \texttt{theorem distanceToSet\_isGLB ... :}\\
  \texttt{\ \ IsGLB (distanceSet M x S) (distanceToSet M x S) := by}\\
  \texttt{\ \ exact Real.isGLB\_sInf}\\
  \texttt{\ \ \ \ (distanceSet\_nonempty M x set\_nonempty)}\\
  \texttt{\ \ \ \ (distanceSet\_bddBelow M x S).}
\end{gathered}
\]
This is the point where Lean uses the completeness of the real numbers. The
metric axioms supply nonnegativity; the real-number library supplies the
existence and correctness of the infimum for nonempty bounded-below subsets of
\(\mathbb{R}\).
\end{remark*}

\begin{remark*}
The next basic theorem says that the distance from a point to a set is bounded
above by the distance from that point to any particular witness in the set:
\[
  \operatorname{dist}(x,A) \le d(x,a)
  \qquad\text{whenever } a \in A.
\]
This direction is important. The distance to the set is the infimum of all
distances from \(x\) to points of \(A\), so it must be less than or equal to
each particular distance in that collection.

In Lean, the theorem is:
\[
\begin{gathered}
  \texttt{theorem distanceToSet\_le\_distance\_to\_point\_of\_mem}\\
  \texttt{\ \ \ \ \{Point : Type u\}}\\
  \texttt{\ \ \ \ (M : MetricSpace Point)}\\
  \texttt{\ \ \ \ (x : Point)}\\
  \texttt{\ \ \ \ \{A : Set Point\}}\\
  \texttt{\ \ \ \ \{a : Point\}}\\
  \texttt{\ \ \ \ (point\_in\_set : a in A) :}\\
  \texttt{\ \ \ \ distanceToSet M x A <= M.distance x a := by}\\
  \texttt{\ \ have A\_nonempty : A.Nonempty := <a, point\_in\_set>}\\
  \texttt{\ \ have distance\_value\_in\_set :}\\
  \texttt{\ \ \ \ M.distance x a in distanceSet M x A := by}\\
  \texttt{\ \ \ \ exact <a, point\_in\_set, rfl>}\\
  \texttt{\ \ have infimum\_property :}\\
  \texttt{\ \ \ \ IsGLB (distanceSet M x A) (distanceToSet M x A) := by}\\
  \texttt{\ \ \ \ exact distanceToSet\_isGLB M x A\_nonempty}\\
  \texttt{\ \ exact infimum\_property.1 distance\_value\_in\_set.}
\end{gathered}
\]
No separate nonemptiness hypothesis is needed here: the assumption
\texttt{point\_in\_set : a in A} already provides a witness that \(A\) is
nonempty.

The proof follows the ordinary infimum argument. Since \(a \in A\), the number
\(d(x,a)\) belongs to the distance set
\[
  \{d(x,y) : y \in A\}.
\]
The theorem \texttt{distanceToSet\_isGLB} says that
\(\operatorname{dist}(x,A)\) is the greatest lower bound of this distance set.
In particular, it is a lower bound. Therefore it is less than or equal to every
member of the distance set, including \(d(x,a)\).

The Lean proof mirrors these sentences. First,
\[
  \texttt{A\_nonempty : A.Nonempty}
\]
is obtained from the witness \(a \in A\). Second,
\[
  \texttt{distance\_value\_in\_set}
\]
records that \texttt{M.distance x a} is one of the values in
\texttt{distanceSet M x A}. Third,
\[
  \texttt{infimum\_property}
\]
applies \texttt{distanceToSet\_isGLB}. Finally,
\[
  \texttt{infimum\_property.1}
\]
is the lower-bound half of the \texttt{IsGLB} statement. Applying it to
\texttt{distance\_value\_in\_set} gives exactly
\[
  \texttt{distanceToSet M x A <= M.distance x a}.
\]

This theorem is often the first useful estimate involving point-to-set
distance. It lets us replace the abstract infimum \(\operatorname{dist}(x,A)\)
with any concrete distance \(d(x,a)\) to a chosen point of \(A\). Some
corollaries of this kind of estimate will be explored later in Volume IV.
\end{remark*}

\begin{remark*}
The next theorem says that the distance from a point to a set vanishes when the
point already belongs to the set:
\[
  x \in A \quad\Longrightarrow\quad \operatorname{dist}(x,A) = 0.
\]
In Lean:
\[
\begin{gathered}
  \texttt{theorem distanceToSet\_eq\_zero\_of\_mem}\\
  \texttt{\ \ \ \ \{Point : Type u\}}\\
  \texttt{\ \ \ \ (M : MetricSpace Point)}\\
  \texttt{\ \ \ \ \{A : Set Point\}}\\
  \texttt{\ \ \ \ \{x : Point\}}\\
  \texttt{\ \ \ \ (point\_in\_set : x in A) :}\\
  \texttt{\ \ \ \ distanceToSet M x A = 0 := by}\\
  \texttt{\ \ ...}
\end{gathered}
\]

The proof is an equality proof by two inequalities. First,
\[
  \operatorname{dist}(x,A) \le d(x,x)
\]
by \texttt{distanceToSet\_le\_distance\_to\_point\_of\_mem}, using the witness
\(x \in A\). Since \(d(x,x)=0\), this gives
\[
  \operatorname{dist}(x,A) \le 0.
\]

Second, every metric distance is nonnegative, so \(0\) is a lower bound for
\[
  \{d(x,y) : y \in A\}.
\]
Since \(\operatorname{dist}(x,A)\) is the greatest lower bound of this set, it
is greater than or equal to every lower bound. In particular,
\[
  0 \le \operatorname{dist}(x,A).
\]
Combining the two inequalities gives
\[
  \operatorname{dist}(x,A)=0.
\]

The Lean proof names these two halves:
\[
\begin{gathered}
  \texttt{have distanceToSet\_nonnegative :}\\
  \texttt{\ \ 0 <= distanceToSet M x A := by}\\
  \texttt{\ \ ...}\\
  \texttt{have distanceToSet\_nonpositive :}\\
  \texttt{\ \ distanceToSet M x A <= 0 := by}\\
  \texttt{\ \ ...}\\
  \texttt{exact le\_antisymm distanceToSet\_nonpositive distanceToSet\_nonnegative.}
\end{gathered}
\]
The final command \texttt{le\_antisymm} is the standard way to prove equality
in an ordered type: prove \(a \le b\) and \(b \le a\).

For the nonnegative half, Lean uses \texttt{distanceToSet\_isGLB}. The lower
bound \(0\) is expressed as
\[
  \texttt{0 in lowerBounds (distanceSet M x A)}.
\]
To prove this, Lean introduces an arbitrary member of the distance set. Since
\texttt{distanceSet M x A} is an image, that member has the form
\texttt{M.distance x y} for some \texttt{y in A}. The metric axiom
\texttt{M.metric.distance\_nonnegative x y} then proves it is at least \(0\).

For the nonpositive half, Lean chains two facts:
\[
\begin{gathered}
  \texttt{distanceToSet M x A <= M.distance x x}\\
  \texttt{M.distance x x = 0.}
\end{gathered}
\]
The first is the previous witness theorem applied to \texttt{point\_in\_set};
the second is \texttt{M.distance\_self x}.
\end{remark*}

\begin{remark*}
A natural follow-on corollary goes in the other direction for a chosen witness
point, but it needs more structure. If \(a \in A\), then one expects
\[
  d(x,a) \le \operatorname{dist}(x,A) + \operatorname{diam}(A).
\]
This says that once \(x\) gets close to the set \(A\), any particular point
\(a \in A\) is at most one diameter of \(A\) farther away.

This is not a direct consequence of
\[
  \operatorname{dist}(x,A) \le d(x,a).
\]
It uses the triangle inequality and the definition of diameter. The hand proof
has the following shape. For any \(y \in A\),
\[
  d(x,a) \le d(x,y) + d(y,a)
\]
by the triangle inequality. Since \(y \in A\) and \(a \in A\), the distance
\(d(y,a)\) is one of the pairwise distances used to define
\(\operatorname{diam}(A)\). Therefore
\[
  d(y,a) \le \operatorname{diam}(A).
\]
So for every \(y \in A\),
\[
  d(x,a) \le d(x,y) + \operatorname{diam}(A).
\]
Passing to the infimum over all \(y \in A\) gives
\[
  d(x,a) \le \operatorname{dist}(x,A) + \operatorname{diam}(A).
\]

In Lean, this theorem will involve both \texttt{distanceToSet} and
\texttt{diameter}. It should therefore not be placed in either
\texttt{SetDistance.lean} or \texttt{SetDiameter.lean} if doing so would create
an artificial dependency. A better location is a third interaction file, such
as
\[
  \texttt{SetDistanceDiameter.lean}.
\]
That file can import both foundational files:
\[
\begin{gathered}
  \texttt{import ...SetDistance}\\
  \texttt{import ...SetDiameter.}
\end{gathered}
\]
A good theorem name for the estimate is:
\[
  \texttt{distance\_to\_point\_le\_distanceToSet\_add\_diameter\_of\_mem}.
\]
The name follows the inequality:
\[
  \texttt{M.distance x a <= distanceToSet M x A + diameter M A},
\]
and the suffix \texttt{of\_mem} records the membership assumption
\texttt{a in A}.
\end{remark*}

\begin{remark*}
The bounded-below fact is likely to be useful elsewhere, but it should live
near \texttt{distanceSet}, not in the core metric-space file. The core metric
file can provide facts about \texttt{M.distance x y}; it should not need to know
about point-to-set constructions. The reusable lemma here is therefore
\[
  \texttt{distanceSet\_bddBelow}.
\]
A theorem about the greatest lower bound should remain separate:
\[
  \texttt{distanceToSet\_isGLB}.
\]
The bounded-below lemma proves only that the set of distances has some lower
bound. The GLB theorem additionally identifies \texttt{distanceToSet M x S} as
the greatest lower bound of that set.
\end{remark*}

\begin{remark*}
This is a typical Lean pattern: definitions, lemmas, and theorems are arranged
so that each later proof can cite the exact earlier fact it needs.

At the metric-space level, the basic reusable fact is about the distance
function itself:
\[
\begin{gathered}
  \texttt{M.distance\_image\_bddBelow x S :}\\
  \texttt{\ \ BddBelow ((fun y : Point => M.distance x y) '' S).}
\end{gathered}
\]
This theorem knows nothing about point-to-set distance. It only says that the
image of any set under the function \(y \mapsto d(x,y)\) is bounded below by
\(0\). The proof uses only the nonnegativity axiom of a metric.

At the set-distance level, we define
\[
  \texttt{distanceSet M x S}
\]
to be exactly that image:
\[
  \texttt{(fun y : Point => M.distance x y) '' S}.
\]
Because of this choice, the bounded-below theorem for distance sets becomes a
short wrapper:
\[
\begin{gathered}
  \texttt{theorem distanceSet\_bddBelow ... :}\\
  \texttt{\ \ BddBelow (distanceSet M x S) := by}\\
  \texttt{\ \ exact M.distance\_image\_bddBelow x S.}
\end{gathered}
\]
No metric proof has to be repeated. The old theorem applies because the new
definition was chosen to expose the same mathematical shape.

Then the final theorem
\[
  \texttt{distanceToSet\_isGLB}
\]
does not need to prove boundedness from scratch. It can cite
\texttt{distanceSet\_bddBelow}. It also cites \texttt{distanceSet\_nonempty},
which comes from the hypothesis \texttt{S.Nonempty}. Once those two inputs are
available, the real-number theorem \texttt{Real.isGLB\_sInf} supplies the
infimum conclusion.

So the dependency chain is:
\[
\begin{array}{c}
  \text{metric nonnegativity}\\
  \Downarrow\\
  \texttt{M.distance\_image\_bddBelow x S}\\
  \Downarrow\\
  \texttt{distanceSet\_bddBelow M x S}\\
  \Downarrow\\
  \texttt{distanceToSet\_isGLB M x set\_nonempty.}
\end{array}
\]
This is why small lemmas are useful in formalization. A lemma isolates a piece
of reasoning that is true for a broad reason. A theorem at a later layer can
then cite the lemma instead of rebuilding that reasoning. The result is closer
to a clean mathematical proof: each line says which previously established fact
is being used.
\end{remark*}

\begin{example*}
On the real line with its usual metric,
\[
  d(x,y) = |x-y|,
\]
the distance from $1$ to the open interval $(2,3)$ is
\[
  \operatorname{dist}(1,(2,3)) = 1.
\]
The set of distances is
\[
  \{ |1-y| : 2 < y < 3 \}.
\]
Since every such $y$ satisfies $y>2$, every distance is greater than $1$:
\[
  |1-y| = y-1 > 1.
\]
But the values can be made arbitrarily close to $1$ by choosing points of the
interval close to $2$ from the right, for example $y = 2+\varepsilon$ with
$0<\varepsilon<1$. Then
\[
  |1-(2+\varepsilon)| = 1+\varepsilon.
\]
Thus $1$ is not a minimum value of the distance set, but it is the greatest
lower bound.

The corresponding Lean test target is:
\[
\begin{gathered}
  \texttt{theorem realMetricSpace\_distanceToSet\_one\_Ioo\_two\_three :}\\
  \texttt{\ \ distanceToSet realMetricSpace (1 : Real) (Set.Ioo (2 : Real) 3) = 1.}
\end{gathered}
\]
To prove this in Lean, one proves that \texttt{1} is an
\texttt{IsGLB} of the set
\[
  \texttt{\{ r : Real | exists y : Real, y in Set.Ioo 2 3 and}
  \quad
  \texttt{r = realMetricSpace.distance 1 y \}}.
\]
The lower-bound half uses $2<y$ to show $1 \le |1-y|$. The greatest-lower-bound
half uses the Archimedean or density properties of the real numbers to find a
point $y \in (2,3)$ whose distance from $1$ is below any proposed lower bound
strictly larger than $1$.
\end{example*}
